\documentclass[11pt]{article}
\usepackage[T1]{fontenc}
\usepackage[utf8]{inputenc}
\usepackage{lmodern}
\usepackage[margin=1in]{geometry}
\usepackage{enumitem}
\usepackage{microtype}
\setlist{nosep,leftmargin=1.4em}
\setlength{\parskip}{2pt}
\setlength{\parindent}{0pt}
\title{\vspace{-3.5em}The Next Step for v16.0: Material Gaps against v11.0\\[2pt]
\large An archaeology review --- TuringAward}
\author{}
\date{2026-07-15}
\begin{document}
\maketitle
\vspace{-3em}

\section{Method}
I inventoried the whole v11.0 corpus and classified each material item against the v16
artifact. \textbf{v11.0 read:} the unified spec body (\texttt{ledger/drafts/sec01--20},
\texttt{appA--I}, via the file set and the signed \texttt{CHANGELOG\_v10.3\_to\_v11.0.md});
\texttt{deferredSettlement.tex} (2129 lines) with its Pareto declaration;
\texttt{market\_data\_mvp\_v1.0/v2.0.tex}; \texttt{data/ledger\_data\_v1.0.tex};
\texttt{valuation/ledger\_valuation\_v1.0.tex}; the convergence and sign-off memos.
\textbf{v16 artifact:} \texttt{Ledger\_Spec\_v16.0/ledger/ledger\_v16\_0.tex} --- the single
consolidated, certified file (110\,pp). It is the only v16 candidate in the repository, so no
panel ruling on artifact identity was needed; the identification was unambiguous and is
recorded here as required. Each v11 item was marked PRESENT (cited in v16), TRANSFORMED
(cited, with what changed), or ABSENT. Every ABSENT item was then subclassified MISSING
(valuable, no evidence of deliberate removal) versus RETIRED, using the v15-era
\texttt{EXCLUSIONS.md} register (E1--E76), the named companions in v16 Ch.~17
(Worked-Examples, SBL-Operations, Track-B, Managed-Account, Simulation), and the out-of-scope
list as the deliberate-retirement trail. RETIRED items are dropped unless their retirement is
argued a mistake; none met that bar. \emph{Finding: v16 is a remarkably complete absorption of
v11.} The vast majority of v11 is PRESENT or TRANSFORMED --- the three-home model, the futures
variation-margin identity, the coordinate vector and coverage invariant, corporate actions and
the market-data operator, close-out netting, the settlement-obligation unit with its
partial-edge split and cum/ex market-claim, virtual ledgers and the TRS, bitemporal time
travel, and an executable-check regime with firing witnesses. Five gaps survived the filter.

\section{Material Gaps}

\textbf{NS-01 --- The settlement--custody reconciliation identity is gone.}
v11's deferred-settlement spec carried an explicit algebraic identity binding the ledger's
internal balances to the external nostro/depot books across the open settlement window:
$\mathrm{nostro}_{\mathrm{ext}} = \mathrm{own}(\mathrm{ccy}) + \sum_{\mathrm{cpty}}
w_{PS}[\cdot] - \mathrm{inflight\_out} + \mathrm{inflight\_in}$, with a depot analogue, both
computable in time bounded by the wallet count rather than the transaction count
(\texttt{deferredSettlement.tex} \S8.1--8.5, invariant DS3; the Conservation Lifting Theorem
\S8.8 makes the identity a corollary). The same identity appears independently in the SBL chapter as the custodian check
$\mathrm{own}+\mathrm{borr}+\mathrm{inflight}=\mathrm{depot}$ (\texttt{sec16}, L466), and
appendix~C lists nostro/vostro and custodian/depot breaks as failures the design detects.
\emph{Absent in v16:} Ch.~11 re-books \texttt{owned} at
instruction and carries the instruction-to-settlement gap in the settlement-obligation unit,
but states no reconciliation identity to the custodian; ``reconcile against the custodian'' is
named at the boundary and never made computable (the closest thing, dual valuation's per-venue
price check, reconciles \emph{value}, not the position/cash balances). This is the one
internal--external check the constitution keeps \emph{in} scope (reconciled against at the
boundary), yet it is the ledger's only defence against the exact drift that motivated the
deferred-settlement work. \emph{Lens:
L2 (transactions and data --- reconciliation).} \emph{Reinstate:} add one subsection to Ch.~11
stating the nostro/depot identity as a projection over \texttt{owned} plus the open
settlement-obligation units per counterparty, with the three-way classify (clean /
expected lead-lag / break) that distinguishes an open-window state from a true break, and one
executable property asserting the identity holds within tolerance on generated settlement
histories. \emph{Cost:} $\approx$1\,pp, no new primitive (it reads existing coordinates and
units). \emph{Risk:} low; the risk of omission is higher --- a settlement book with no
computable custody check re-admits the reconciliation failure the framework exists to abolish.

\textbf{NS-02 --- No reference currency, FX translation, or currency minor-unit variation.}
Four v11 documents carry a currency dimension: the data layer registers currencies with
per-currency \texttt{minor\_units} (USD 2, JPY 0) (\texttt{ledger\_data\_v1.0.tex} $L_1$);
the market-data MVP works a bitemporal EUR/USD restatement to the pip; valuation calibrates FX
state (3--5 latent factors); and deferred settlement models a cross-currency trade as a parent
FX obligation with two child legs, the open legs a \emph{quantified Herstatt exposure}
(\texttt{deferredSettlement.tex} \S7.4, DS10). \emph{Absent in v16:} the specification assumes a
single implicit currency throughout --- ``exact integers in minor units,'' ``exact in cents''
--- with no reference-currency NAV, no FX-rate observation feeding a translation operator, and
no cross-currency settlement leg or Herstatt visibility. A book valued and settled in one
currency is not the trading book in scope. The ``cents'' assumption is itself a latent defect:
rounding a JPY figure to two decimals is a hundred-fold error. \emph{Lens: L9 (numerical
computation --- currency arithmetic and minor units).} \emph{Reinstate:} state that money is an
integer count at the currency's declared minor-unit scale (a datum of the currency unit, not a
global constant); admit the FX rate as an ordinary recorded observation; make reference-currency
NAV an FX-translation projection over the same composition; and carry a cross-currency
settlement obligation as two legs whose open state projects as Herstatt exposure. \emph{Cost:}
$\approx$1--1.5\,pp across Ch.~7 (valuation) and Ch.~11 (settlement). \emph{Risk:} medium ---
FX translation touches the valuation and settlement projections, but adds no new primitive (an
FX rate is an observation, a translation is pricing).

\textbf{NS-03 --- The client-isolation (segregation) theorem, and its false-theorem warning,
were dropped.}
The v10.3$\to$v11.0 changelog records a correctness-critical correction: v10.3's ``conservation
alone enforces segregation by algebra'' was a \emph{false theorem} --- it permits a cross-client
move --- replaced in v11 (\texttt{sec09}, \texttt{thm:seg}, L240) by segregation as a theorem holding
only under $\mathrm{CONS}\wedge\mathrm{LOC}\wedge\mathrm{C4}$ (conservation, locality, and the
capability-read condition); the same ``conservation alone is insufficient'' warning drives the
singleton-bilateral mandate invariant (L44). The v15 exclusions register (E31) explicitly
promised ``the segregation theorem itself is carried.'' \emph{Absent in v16:} client isolation
\emph{within} a ledger rests on a single sentence --- ``authorisation is checked against the
manager'' (Ch.~3, Wallets) --- with no theorem, no locality/capability precondition, and no
record of the warning that conservation alone is \emph{insufficient}. (v16 does carry the
\emph{cross-ledger} realm isolation, v11's P7 --- ``the bridge is an observation, never a
move'' (Ch.~10) --- but that is a different property: it isolates the virtual ledger from the
true one, not one client's wallets from another's inside one ledger.) The gap is not
decorative: it is the property under which one manager cannot move another beneficial owner's
assets. A register that claims an item is carried while the artifact omits
it is the strongest possible evidence of a MISSING, not a RETIRED, item. \emph{Lens: L1
(concurrency and isolation --- and L8 capability; the safety property a single writer does not
by itself supply).} \emph{Reinstate:} state the isolation theorem in Ch.~6 or Ch.~4 ---
authorised balance change is confined to wallets the acting manager holds a capability for, so
conservation \emph{plus} the capability-read gate \emph{plus} move-locality makes a cross-client
move unrepresentable --- with an executable property that an unauthorised cross-wallet move is
refused. \emph{Cost:} $<$0.5\,pp. \emph{Risk:} low; it names a guarantee the door already
enforces but the prose leaves as an aside.

\textbf{NS-04 --- The recursive partial-settlement split has no termination or
fragmentation bound.}
v11 bounded partial-fill fragmentation: a partial cascade terminates within $D_{\max}=2$ else
the residual goes \texttt{Failed}$\to$\texttt{Compensated}, with per-step conservation and
monotonic advance (\texttt{deferredSettlement.tex} \S6.3, invariants DS11a/DS11b).
\emph{Diminished in v16:} Ch.~11's partial edge splits a settlement-obligation unit ``again by
the same edge, recursively,'' with no depth or fragment-count bound stated. It terminates only
because quantity is a finite integer, so a $Q$-unit instruction can fragment into up to $Q$
live obligation units, each entering the obligation-liveness queue --- an operational and
liveness cost the v11 bound existed to cap. \emph{Lens: L3 (specification and verification ---
liveness and a variant function for termination).} \emph{Reinstate:} add the fragment bound and
its variant to the partial edge (residual strictly decreases; beyond the declared depth the
residual is failed-and-compensated rather than re-split), and property-test that no generated
settlement history exceeds it. \emph{Cost:} a few lines in Ch.~11 and one property in Ch.~15.
\emph{Risk:} low; it strengthens an existing edge and cannot weaken a proved property.

\textbf{NS-05 --- The economic-recomputation oracle does not pin the contract-code version.}
v16's economic check re-runs the map on the recorded event and compares:
\texttt{contract e (ledgerBefore e) === txRecorded} (Ch.~15, \texttt{prop\_recompute}). This is
sound only against the contract as it stood when the event was admitted. v11's data layer pinned
this explicitly: a \texttt{VersionPin} with a \emph{contract} axis alongside schema, model, and
canonicalisation, so a recomputation years later is deterministic
(\texttt{ledger\_data\_v1.0.tex} \S5.2). \emph{Absent in v16:} ProductTerms are versioned and
contracts are called ``versioned functions,'' but nothing records \emph{which} contract-code
version a recomputation must use; an amended contract silently re-derives a different
transaction and the oracle either false-fails or false-passes. Replay is safe (it re-reads
transactions, never re-executes contracts), but the \emph{economic} oracle re-executes and is
therefore under-specified. \emph{Lens: L3 (verification --- reproducibility of the check
itself).} \emph{Reinstate:} record the contract-code version on the transaction (or bind it to
the ProductTerms version already recorded) and state that \texttt{prop\_recompute} reads the
contract version in force at that log position. \emph{Cost:} one sentence in Ch.~15 plus a field
note in Ch.~5. \emph{Risk:} low; it closes a determinism hole in a check the specification
already relies on.

\section{Declaration}
\textbf{Lenses applied:} L1 (isolation --- NS-03), L2 (reconciliation --- NS-01), L3
(verification/liveness --- NS-04, NS-05), L9 (numerical/currency --- NS-02); L8 (capability)
considered jointly with NS-03. \textbf{Lenses dismissed:} L4 (algorithms/complexity --- no
hidden hard subproblem; the fold is linear); L5 (abstraction/language --- v16's type discipline
is sound and idiomatic); L6 (systems architecture --- conceptual integrity is intact, one
document, one writer); L7 (networks --- no transport is specified; delivery is the untrusted
machines' concern, already handled); L10 (learning --- no learned component is load-bearing for
any safety property, by design). \textbf{Confidence:} MEDIUM-HIGH. The v16 read was complete and
the v11 companion inventories were thorough; the residual unknown is whether NS-01 (the
custody reconciliation identity) is judged in-scope or an external-authority matter the owner
intends to leave at the perimeter --- I read it as in-scope because the constitution keeps
boundary reconciliation in scope, but that is the one call an owner could decide the other way.
\textbf{Decision log:} DL-01 $\mid$ archaeology $\mid$ which file is the authoritative v16
artifact? $\mid$ single certified consolidated file present, no rival $\mid$ ruling:
\texttt{ledger\_v16\_0.tex}, unambiguous $\mid$ no panel convened; no ambiguity met the
threshold for a vote.
\end{document}
